\SourceLeaf{322}
\[
\begin{array}{rcll@{\qquad}rcll}
2 & \text{um} & 180^0+A & \text{rechts von }1 & 3 & \text{um} & B & \text{rechts von }1\\[2pt]
 & \text{,,} & A & \text{,,} & & \text{,,} & 180^0+B & \text{,,}\\[2pt]
 & \text{,,} & 180^0+A & \text{,,} & & \text{,,} & 180^0+B & \text{,,}\;.
\end{array}
\]

Das Ergebnis selbst stimmt aber immer nur mit einem dieser vier Datensätze überein, und die anderen drei sind \textls[120]{physisch unmöglich}.

Der metaphysische Grund dieser Erscheinung ist, dass man von der beobachteten Richtung nur den Umstand benutzt, dass die \textls[120]{geraden Linien} bestimmte Winkel miteinander bilden. Dabei lässt man diese geraden Linien nach \textls[120]{beiden} Seiten in indefinitum verlaufen, während das Licht nur in einem bestimmten Sinn fortschreitet, sodass man die Fälle ausschließen muss, in denen der Gegenstand in der entgegengesetzten Richtung läge.

Zur Vollständigkeit gehört also, die Lösung der Aufgabe, schon vor der Rechnung unmittelbar an den Daten zu erkennen, ob sie vereinbar sind. Das lässt sich sehr elegant erledigen, doch möchte ich Ihnen bei der eigenen Herleitung nicht vorgreifen. \mbox{. . . . . . .}



\centerline{\textsc{Gauss} an \textsc{Gerling}. Göttingen, 14. Januar 1842.}

\indent \mbox{. . . . . . . .} Das Kriterium für die \textsc{Pothenot}sche Aufgabe teile ich Ihnen gern mit, wünsche jedoch, dass Sie es für sich behalten. Der Grund dafür ist, dass ich mir vorbehalte, das Theorem, mit dem es zusammenhängt, bei passender Gelegenheit selbst zu behandeln, weniger wegen der Eleganz des Theorems an sich als wegen der Eleganz, die sich hier aus der Anwendung komplexer Größen ergibt. Es sollte also insbesondere eine Gelegenheit sein, bei der ich mehr über den Gebrauch komplexer Größen sagen kann. Würde das Theorem vorher von anderer Seite veröffentlicht, so hätte ich nachher keine Lust mehr, selbst darauf zurückzukommen. Nun aber zur Sache.

Man kann das Kriterium in vielerlei Formen ausdrücken. Am elegantesten ist es wohl, alle drei gegebenen Punkte darin symmetrisch auftreten zu lassen. Ich bezeichne also die drei gegebenen Punkte mit 1, 2, 3;\space%
